\section{Implementation}
\label{sec:implementation}
We implemented \textsc{ToFU} as an 8,400-SLOC Python framework (Table~\ref{tab:implementation-sloc}, excluding tests and libraries), running on Python 3.12, Atheris v2.3, LangChain Core v0.3, the official \texttt{mcp} SDK, and Node.js v20+.

\begin{table}[t]
\centering
\caption{Codebase breakdown by major component.}
\label{tab:implementation-sloc}
\small
\resizebox{\columnwidth}{!}{%
\begin{tabular}{llr}
\toprule
\textbf{Component} & \textbf{Modules \& Descriptions} & \textbf{SLOC} \\
\midrule
Sanitizers \& Oracles & \texttt{src/sanitizers/} (21 CWE introspection hooks \& checks) & 3,580 \\
Target Runners & \texttt{src/runners/} (In-process Python \& Node.js/MCP runners) & 2,700 \\
Fuzzing Engine & \texttt{src/fuzzer.py}, \texttt{src/input\_parser.py} (Schema-aware mutation) & 1,650 \\
Support \& Telemetry & \texttt{src/metrics.py}, \texttt{src/reporters.py}, \texttt{src/llm.py} (Seeding \& logs) & 470 \\
\midrule
\textbf{Total Core Code} & \textit{Excludes tests, prompts, and third-party tools} & \textbf{8,400} \\
\bottomrule
\end{tabular}%
}
\end{table}

\noindent\textbf{Sanitizer Specialization \& LLM Seeding.}
\textsc{ToFU} prompts \mbox{MiniMax M3} ($T=0.7$, up to 4 retries, fixed prompt in Appendix~\ref{sec:appendix-prompts}) with the target tool's metadata, JSON schema, and attack vectors from an internal seed bank to produce $N=6\text{--}12$ structured seeds. To synthesize valid binary documents, \textsc{ToFU} exposes native Python tool bindings (\texttt{reportlab} for PDF, \texttt{python-docx} for DOCX, \texttt{openpyxl} for XLSX, \texttt{zipfile} for archives); the LLM invokes these tools via standard function calling, materializing files to disk and referencing them as \texttt{predefined\_seed:<path>}. A deterministic fallback populates boundary numbers ($0, -1$), boolean toggles, and malformed types if the LLM fails ($<2\%$ of queries). Concurrently, \textsc{ToFU} matches schema properties against a 21-CWE corpus (\texttt{src/sanitizers/}), compiling executable Python verification predicates cached per schema hash.

\noindent\textbf{Fuzzing Engine \& Runtime Interception.}
\textsc{ToFU} recursively decodes libFuzzer-mutated byte streams into schema-valid arguments using \texttt{atheris.FuzzedDataProvider} (\texttt{FuzzInputParser}), mapping strings (\texttt{ConsumeUnicodeNoSurrogates}), numbers, arrays, objects, and deliberate type-mismatches. Feedback is measured as unique Python bytecode branch transitions via \texttt{sys.settrace} (and Atheris C-hooks), strictly filtering out framework internals (\texttt{src/fuzzer.py}, \texttt{src/sanitizers/}). To extract libFuzzer feature updates without disk I/O or terminal contention, C-level \texttt{stderr} is duplicated (\texttt{dup2}) into an IPC pipe drained asynchronously by a worker thread. Dual oracles evaluate each run: (1)~\emph{Runtime Introspection}, where \texttt{ExitStack} wraps execution sinks (\texttt{open}, \texttt{subprocess}, \texttt{socket}, \texttt{asyncio.get\_running\_loop}) to evaluate path containment, argument injection, and async loop freezing; and (2)~\emph{Black-Box Differential}, checking schema contracts and silent failures (\texttt{failed\_init\_non\_exit}).

\noindent\textbf{Target Execution \& Telemetry.}
\textsc{ToFU} supports in-process execution for Python tools (\texttt{PythonRunner}) and stdio subprocess workers for JS/TS MCP servers (\texttt{JSRunner}). Each iteration runs under a 5-second watchdog in an isolated temporary directory (\texttt{mkdtemp}) with environment variable restoration. Crashes are deduplicated by hashing the top-3 normalized target stack frames grouped by triggered sanitizer ID. Every trial logs execution latency, input values, triggered sanitizers, return status, edge coverage, and crash signatures.